\subsection{贝尔空间}

定义 3. 拓扑空间 X 称为贝尔空间，如果它满足下列两个等价条件之一：

（EB） \(\mathbf{X}\) 中稠密开集的每个可数交在 \(\mathbf{X}\) 中稠密。

（EB'）X 中没有内点的闭集的每个可数并在 X 中没有内点。

公理 (EB) 还可以表述成另外两个等价形式：

（EB''）X 中每个非空开集都不是贫集。

因为一个集是贫集，当且仅当它含于可数多个没有内点的闭集的并。

（EB''\,``）X 中贫集的余集在 X 中稠密。

这表示贫集不能包含非空开集，因而等价于 \((EB'')\) 。

命题 3. 贝尔空间 X 的每个非空开子空间 Y 都是贝尔空间。

这由 (EB'') 得出，因为 Y 中每个开集（相应地，贫集）在 X 中是开集（相应地，贫集）。

由命题 3 推出，贝尔空间的每一点都有一个基本邻域系，其中每个邻域都是贝尔空间。反之：

命题 4. 若拓扑空间 X 的每一点都有一个邻域是贝尔空间，则 X 是贝尔空间。

设 A 是 X 中的非空开集，设 \(x \subset A\) ，并设 V 是 x 的开邻域且是贝尔空间。若 A 在 X 中是贫集，则 \(V \cap A\) 在 V 中是贫集且在 V 中开，这与假设相矛盾。

命题 5. 在贝尔空间 X 中，贫集的余集是贝尔空间。

设 A 是 X 中的贫集；则它在 X 中的余集 Y = CA 在 X 中稠密。设 B 是相对于 Y 的贫集；B 相对于 X 也是贫集，故 A ∪ B 相对于 X 是贫集。于是 A ∪ B 在 X 中的余集（它也是 B 在 Y 中的余集）在 X 中稠密，因而在 Y 中更稠密。故 Y 是贝尔空间。

定理 1（贝尔）。(i) 每个局部紧空间 X 都是贝尔空间。(ii) 每个这样的拓扑空间 X 都是贝尔空间：在 X 上存在一个与 X 的拓扑相容的度量，且它使 X 具有完备度量空间的结构。

我们来证明，在每种情形公理 (EB) 都满足。设 \((\mathbf{A}_n)\) 是 \(\mathbf{X}\) 中稠密开集的序列，并设 \(\mathbf{G}\) 是任一非空开集。于是我们可以归纳地定义一个非空开集的序列 \((\mathbf{G}_n)\) ，使得 \(\mathbf{G}_1 = \mathbf{G}\) 且 \(\overline{\mathbf{G}}_{n+1} \subset \mathbf{G}_n \cap \mathbf{A}_n\) ：因为按假设 \(\mathbf{G}_n\) 非空，故 \(\mathbf{G}_n \cap \mathbf{A}_n\) 是非空开集；而由于在所考虑的两种情形中 \(\mathbf{X}\) 都是正则的，故存在非空开集 \(\mathbf{G}_{n+1}\) 使 \(\overline{\mathbf{G}}_{n+1} \subset \mathbf{G}_n \cap \mathbf{A}_n\) 。于是集 \(\mathbf{G} \cap \bigcap_{n=1}^{\infty} \mathbf{A}_n\) 包含诸集 \(G_{n}\) 的交，而后者等于诸集 \(\overline{G}_{n}\) 的交；因此只需证明 \(\bigcap_{n=1}^{\infty}\overline{G}_{n}\neq\varnothing\) 。现若 X 局部紧，可设 \(\overline{G}_{2}\) 是紧的；在紧空间 \(\overline{G}_{2}\) 中，诸 \(\overline{G}_{n}(n\geqslant2)\) 构成一个非空闭集的递减序列，故由公理 (C'\,') 它们的交非空（参看第 I 章，§9，第 1 号，定义 1{]}。若 X 是完备度量空间（关于一个与其拓扑相容的度量），则可设 \(\overline{G}_{n}\) 已选得使其直径（关于该度量）当 n 趋于 \(+\infty\) 时趋于 0 ；于是诸 \(\overline{G}_{n}\) 构成一个柯西滤子基，它收敛于一点 x ，而 x 必属于诸集 \(\overline{G}_{n}\) 的交。证毕。

注。存在不属于这两类中任何一类的贝尔空间，特别有既不可度量化又非局部紧的贝尔空间（习题 16）；也存在不具有与其拓扑相容的完备度量空间结构的可度量化贝尔空间（习题 14）。

